\gddchapter{Testing sequence}
The prototype can be played with both voice and keyboard and the navigation modes can be swapped with a button press. 

\section{Baseline A - Keyboard navigation}
At first the participant will spend 10 minutes by playing the game using keyboard navigation.
This allows the player to anchor in the typical experience and serves as a comparative point in the A/B testing.

\subsection{Observation log}
\begin{table}[h]
\centering
\begin{tabular}{|p{5cm}|p{10cm}|}
\hline
\textbf{Field} & \textbf{Response} \\ \hline
\textbf{Physical comfort} How is the participant feeling & The player looks comfortable without any signs of distress \\[1.5cm] \hline
\textbf{Immersion markers} (flow \& frustration) & Participant seems to be moderately inside of the game, pulling out occasionally to make remarks or comments to the researcher\\[1.5cm] \hline
\end{tabular}
\end{table}

\section{Novel version B - Voice navigation}
After 10 minutes the researcher switches the navigation to voice input mode by using the keyboard. The prototype is not restarted,
instead the player is instructed to continue his playthrough using voice. The decision not to restart is placed to reduce any frustration that the player
might feel if they have to restart and repeat the actions that they just did.

\subsection{Observation log}
\begin{table}[H]
\centering
\begin{tabular}{|p{5cm}|p{10cm}|}
\hline
\textbf{Field} & \textbf{Response} \\ \hline
\textbf{Physical comfort} How is the participant feeling & Player looks physically comfortable \\[1.5cm] \hline
\textbf{Natural language observation} (is the player using long sentences or simple commands? What type of language do they use? Are they speaking first or 3rd person?) & Participant uses simple commands "Go to X", "Pick up Y". They use only 2nd person\\[1.5cm] \hline
\textbf{Linguistic guessing} (is the particpant trying to guess the possible commands by speaking to the game?) & The participant uses only the commands that were already exposed to them in the baseline version\\[1.5cm] \hline
\textbf{Processing time user reaction} (Does the wait time break the suspension of disbelief)& The disbelief was impacted mostly during a total faliure of the voice recognition pipeline that required a restart of the whole back-end infrastructure.\\[1.5cm] \hline
\textbf{Voice processing pipeline edge cases} (System edge cases eg. two vs to)& As mentioned above the voice interaction pipeline fully stopped working at some point. This was the first ovserved faliure of this scope\\[1.5cm] \hline
\end{tabular}
\end{table}

\subsection{Sticky moments}
The moments that stood out during testing marked with timestamps.
\begin{table}[h]
\centering
\begin{tabular}{|p{3cm}|p{10cm}|}
\hline
\textbf{Timestamp} & \textbf{Log} \\ \hline
03:05 & Start of the playtest \\ \hline 
05:00 & Item usage UX \\ \hline 
07:10 & Fishing store bug \\ \hline 
09:10 & Not said but stairs were confusing \\ \hline 
10:00 & Frustration - I cant't use item in Korean store \\ \hline 
12:00 & Lounge VUI fail \\ \hline 
13:00 & Place before \\ \hline 
16:00 & WhisperCPP broke \\ \hline 
Recording 2: 06:40 & Empty guns bug \\ \hline 



\end{tabular}
\end{table}

The last event on the above table has the timestamp of 06:40 since the recording falied and had to be restarted.


\section{Alignment check}
How much was the researcher participating in the gameplay experience? 

The researcher did his best not to participate in the gameplay experience but he did have to intervene when game breaking bugs arose.



